% Data roles in ensemble development and final evaluation.
% Schematic only; no dataset sizes or empirical quantities are encoded.
% Canvas: 164 mm x 82 mm. Dependencies: project palette and TikZ.
\begin{tikzpicture}[
  x=1mm,y=1mm,
  font=\footnotesize,
  text=BookInk,
  source/.style={
    rectangle,
    draw=FigureInput,
    fill=FigureInput!8,
    line width=0.7pt,
    minimum width=43mm,
    minimum height=9mm,
    text width=39mm,
    align=center,
    inner sep=1mm
  },
  condition/.style={
    rectangle,
    draw=FigureLoss,
    fill=FigureLoss!9,
    line width=0.7pt,
    minimum width=52mm,
    minimum height=9mm,
    text width=48mm,
    align=center,
    inner sep=1mm
  },
  purpose/.style={
    rectangle,
    draw=FigureOutput,
    fill=FigureOutput!6,
    line width=0.8pt,
    minimum width=45mm,
    minimum height=9mm,
    text width=41mm,
    align=center,
    inner sep=1mm
  },
  final/.style={
    source,
    line width=0.9pt
  },
  final purpose/.style={
    purpose,
    line width=0.9pt
  },
  flow/.style={
    -{Latex[length=2.2mm,width=1.5mm]},
    draw=BookMuted,
    line width=0.8pt
  }
]
  \path[use as bounding box] (0,0) rectangle (164,82);

  \node[font=\heiti\small,text=FigureInput] at (25,78) {预测或数据来源};
  \node[font=\heiti\small,text=FigureLoss] at (82,78) {保持隔离的条件};
  \node[font=\heiti\small,text=FigureOutput] at (139,78) {在开发流程中的用途};

  \node[source] (oob-source) at (25,65) {袋外预测};
  \node[condition] (oob-condition) at (82,65)
    {成员$m$未使用样本$i$训练};
  \node[purpose] (oob-purpose) at (139,65)
    {Bagging内部诊断与构造};

  \node[source] (oof-source) at (25,51) {折外预测};
  \node[condition] (oof-condition) at (82,51)
    {折内模型未使用当前留出折};
  \node[purpose] (oof-purpose) at (139,51)
    {训练Stacking元学习器};

  \node[source] (competence-source) at (25,37) {能力估计预测};
  \node[condition] (competence-condition) at (82,37)
    {能力数据不参与成员训练};
  \node[purpose] (competence-purpose) at (139,37)
    {估计局部能力并路由};

  \node[source] (validation-source) at (25,23) {验证预测};
  \node[condition] (validation-condition) at (82,23)
    {与相应候选的拟合路径隔离};
  \node[purpose] (validation-purpose) at (139,23)
    {调参、筛选、校准与早停};

  \draw[BookRule,line width=0.8pt] (1,14)--(163,14);
  \node[fill=white,inner sep=1mm,font=\heiti\small,text=BookInk] at (82,14)
    {冻结全部训练、选择、校准与部署规则};

  \node[final] (test-source) at (25,6) {测试预测};
  \node[condition,line width=0.9pt]
    (test-condition) at (82,6)
    {测试样本未进入任何开发决策};
  \node[final purpose] (test-purpose) at (139,6)
    {评价最终完整流程};

  \foreach \prefix in {oob,oof,competence,validation,test} {
    \draw[flow] (\prefix-source)--(\prefix-condition);
    \draw[flow] (\prefix-condition)--(\prefix-purpose);
  }
\end{tikzpicture}
